\section{Busca de Custo Mínimo e Busca A*}
\label{sec:busca-informada}

Nesta etapa foram implementados dois algoritmos que utilizam o custo das
transições para orientar a exploração do espaço de estados: a Busca de Custo
Mínimo e a Busca A*. Como o objetivo do problema da Ponte e Tocha é minimizar
o tempo total necessário para que todas as pessoas atravessem a ponte, o custo
acumulado até cada estado passa a ser uma informação essencial para a busca.

\subsection{Função de custo acumulado}

Em cada travessia, uma ou duas pessoas podem atravessar a ponte. O custo da
ação corresponde ao tempo da pessoa mais lenta entre aquelas envolvidas no
movimento. Dessa forma, para um conjunto $M$ de pessoas que atravessam, o
custo da transição é definido por

\begin{equation}
	c(M)=\max_{p\in M} t(p),
\end{equation}

em que $t(p)$ representa o tempo necessário para a pessoa $p$ atravessar a
ponte.

A função $g(n)$ representa o custo total acumulado desde o estado inicial até
o estado $n$. Para um sucessor $n'$ de $n$, tem-se

\begin{equation}
	g(n') = g(n) + c(n,n').
\end{equation}

Assim, $g(n)$ corresponde ao tempo efetivamente gasto para alcançar determinada
configuração do problema.

\subsection{Busca de Custo Mínimo}

A Busca de Custo Mínimo seleciona o caminho da fronteira com menor custo
acumulado $g(n)$, implementando a fronteira como uma fila de prioridade por
meio da estrutura \texttt{heapq} da biblioteca padrão do Python.

Na implementação, cada elemento da fronteira armazena o custo acumulado, um
valor auxiliar de desempate, o estado correspondente e o caminho percorrido
até ele. Dessa forma, o elemento retirado da fila é sempre aquele associado ao
menor valor de $g(n)$.

Também foi utilizado um dicionário denominado \texttt{melhor\_custo}, que
armazena o menor custo conhecido para alcançar cada estado. Quando um sucessor
é gerado, calcula-se

\begin{equation}
	g(n') = g(n) + c(n,n').
\end{equation}

Caso esse novo valor seja menor que o custo anteriormente registrado para o
mesmo estado, o dicionário é atualizado e o novo caminho é inserido na fila de
prioridade.

Entradas antigas da fila, correspondentes a caminhos que já foram superados
por alternativas de menor custo, são descartadas durante a execução.

Como os custos das travessias são não negativos, quando o estado objetivo é
retirado da fila de prioridade, o custo associado a ele corresponde ao menor
custo possível para alcançar a solução.

\subsection{Função heurística}

Para a Busca A* foi definida uma função heurística baseada no maior tempo de
travessia entre as pessoas que ainda permanecem no lado de origem da ponte.

A função utilizada é

\begin{equation}
	h(n)=
	\begin{cases}
		0, & \text{se a origem estiver vazia}, \\[2mm]
		\displaystyle\max_{p\in \mathrm{origem}(n)} t(p),
		   & \text{caso contrário}.
	\end{cases}
\end{equation}

A ideia é considerar o tempo da pessoa mais lenta que ainda precisa sair do
lado inicial. Como todos os tempos de travessia são positivos, tem-se também
$h(n)\geq 0$.

Essa propriedade é compatível com a interpretação de $h(n)$ como uma
estimativa do custo restante. Entretanto, a condição essencial para garantir
que a heurística não comprometa a optimalidade do A* é a sua admissibilidade,
isto é, que ela nunca superestime o custo real mínimo até o objetivo.

A heurística é admissível porque a pessoa mais lenta que ainda se encontra na
origem obrigatoriamente precisará atravessar a ponte pelo menos uma vez antes
que o estado objetivo seja alcançado. Portanto, o custo real restante nunca
poderá ser menor que o tempo de travessia dessa pessoa.

Assim,

\begin{equation}
	h(n) \leq h^*(n),
\end{equation}

em que $h^*(n)$ representa o custo real mínimo restante entre o estado $n$ e
o objetivo.

A função também é consistente, isto é, para uma transição entre $n$ e $n'$,

\begin{equation}
	h(n) \leq c(n,n') + h(n').
\end{equation}

Em um movimento de ida, o custo da travessia é suficiente para cobrir uma
eventual redução da estimativa. Em um movimento de volta, pessoas são
adicionadas novamente ao lado de origem e, portanto, o maior tempo presente
nesse lado não diminui.

\subsection{Integração da heurística ao espaço de estados}

Para que a Busca A* utilize corretamente a função heurística, cada estado deve
armazenar uma estimativa correspondente à sua própria configuração.

No estado inicial, a heurística é calculada logo após a criação do estado. Nos
demais casos, primeiro é construído o estado sucessor, considerando as pessoas
movimentadas e a nova posição da tocha. Somente depois é calculado o valor da
heurística correspondente a esse novo estado.

Dessa forma, evita-se utilizar a estimativa de um estado anterior como se ela
pertencesse ao seu sucessor.

A função heurística padrão do projeto foi mantida com valor zero. Assim, o
espaço de estados pode continuar sendo utilizado por algoritmos que não
necessitam de informação heurística.

\subsection{Busca A*}

A Busca A* utiliza simultaneamente o custo acumulado e uma estimativa do custo
restante. Para cada estado é calculado

\begin{equation}
	f(n)=g(n)+h(n),
\end{equation}

em que $g(n)$ representa o custo real já gasto para alcançar o estado, enquanto
$h(n)$ representa uma estimativa do custo restante até o objetivo.

A fronteira do A*, também implementada como fila de prioridade, utiliza
$f(n)$ para determinar qual estado será analisado primeiro. Assim,
diferentemente da Busca de Custo Mínimo, que utiliza somente o custo
acumulado, o A* também considera uma estimativa de quão próximo
o estado está do objetivo.

Na implementação, cada entrada da fila armazena o valor de prioridade, um
valor auxiliar de desempate, o custo acumulado $g(n)$, o estado e o caminho
percorrido.

O dicionário \texttt{melhor\_custo} continua armazenando apenas os valores de
$g(n)$, pois esses valores representam os custos reais conhecidos. Caso um
caminho de menor custo seja posteriormente encontrado para uma configuração já
conhecida, o estado pode ser inserido novamente na fila.

Como o objeto que representa um estado também contém o valor da heurística,
foi utilizada como identificação lógica da configuração a tupla
$(\text{origem},\text{destino},\text{posição da tocha})$.

Essa escolha evita que duas configurações equivalentes sejam interpretadas
como estados distintos apenas por apresentarem valores diferentes no campo da
heurística.

\subsection{Testes e validação da implementação}

Foram utilizados diferentes conjuntos de tempos de travessia para verificar o
funcionamento dos algoritmos. Os custos obtidos pela Busca de Custo Mínimo e
pela Busca A* foram comparados, verificando-se que ambos encontraram os mesmos
custos ótimos.

Os resultados utilizados durante a validação são apresentados na
Tabela~\ref{tab:testes-tarefa3}.

\begin{table}[H]
	\centering
	\caption{Custos encontrados nos testes da Tarefa 3.}
	\label{tab:testes-tarefa3}
	\begin{tabular}{ccc}
		\toprule
		Tempos de travessia & Busca de Custo Mínimo & Busca A* \\
		\midrule
		$[1]$               & 1                     & 1        \\
		$[1,2]$             & 2                     & 2        \\
		$[1,2,5]$           & 8                     & 8        \\
		$[1,2,5,10]$        & 17                    & 17       \\
		$[2,2,2]$           & 6                     & 6        \\
		\bottomrule
	\end{tabular}
\end{table}

Para a instância clássica do problema,

\begin{equation}
	[1,2,5,10],
\end{equation}

os dois algoritmos encontraram uma solução de custo total igual a 17 \text{ minutos}.

Além da comparação dos custos finais, foram verificados os caminhos retornados
pelos algoritmos. Os testes confirmam que a posição da tocha é alternada após
cada movimento, que são movimentadas uma ou duas pessoas por travessia e que a
soma dos custos das transições coincide com o custo total informado pela
busca.

Também foi testada a Busca A* utilizando uma heurística nula, $h(n)=0$.

Nesse caso,

\begin{equation}
	f(n)=g(n),
\end{equation}

fazendo com que o critério de prioridade do A* se torne equivalente ao da
Busca de Custo Mínimo. Os custos obtidos permaneceram iguais, servindo como
uma verificação adicional da implementação.

Os testes também verificaram o valor da heurística no estado inicial, no estado
objetivo e nos estados sucessores, além da propriedade de consistência nas
transições consideradas.

\subsection{Considerações sobre os métodos}

A principal diferença entre os dois algoritmos está no critério utilizado para
ordenar a fronteira.
Na Busca de Custo Mínimo, $\text{prioridade}(n)=g(n)$,
enquanto na Busca A*, $\text{prioridade}(n)=g(n)+h(n)$.

Assim, a Busca de Custo Mínimo considera somente o custo já realizado, enquanto
a Busca A* utiliza também informações específicas do problema para estimar o
custo que ainda falta.

Como a heurística escolhida é admissível, a Busca A* preserva a garantia de
encontrar uma solução ótima. Ao mesmo tempo, a informação heurística pode
orientar a exploração para regiões mais promissoras do espaço de estados.

A comparação quantitativa entre os quatro algoritmos utilizados no trabalho,
considerando custo da solução, número de nós expandidos e tempo de
processamento, será apresentada posteriormente na seção destinada à Tarefa 4.
